% =====================================================================
%  Aufgabe 1 – Vorlage für ein Kapitel. Tipp: ein Satz pro Zeile,
%  dann zeigt Git Änderungen satzweise und Konflikte werden seltener.
% =====================================================================
\chapter{Aufgabe 1: Titel der Aufgabe}
\label{chap:aufgabe1}

\section{Einleitung}
\label{sec:a1-einleitung}

Hier beginnt der Text zu Aufgabe~1.
Wissenschaftliche Quellen werden im Fließtext zitiert \parencite{wilkinson2016}, ebenso Internetquellen \parencite{github-docs-repos}.
Mit \textcite{wilkinson2016} wird der Name in den Satz eingebunden.

\section{Bilder}

\Cref{fig:a1-diagramm} zeigt ein selbst erstelltes Diagramm, \cref{fig:a1-fremdbild} ein übernommenes Bild.

\begin{figure}[htbp]
  \centering
  \includegraphics[width=0.6\linewidth]{beispiel-diagramm}
  \caption{Selbst erstelltes Beispieldiagramm.}
  \label{fig:a1-diagramm}
\end{figure}

\begin{figure}[htbp]
  \centering
  % Für ein fremdes Bild hier die echte Datei einbinden, z.\,B.:
  % \includegraphics[width=0.6\linewidth]{fremdes-bild}
  \fbox{\parbox[c][3cm][c]{0.6\linewidth}{\centering Platzhalter für ein fremdes Bild}}
  \caption{Nicht selbst erstelltes Bild mit Quellenangabe. \bildquelle{wilkinson2016}}
  \label{fig:a1-fremdbild}
\end{figure}

\section{Tabellen}

Die Ergebnisse sind in \cref{tab:a1-ergebnisse} zusammengefasst.
Bei Tabellen steht die Überschrift oberhalb.

\begin{table}[htbp]
  \centering
  \caption{Beispieltabelle mit Tabellenüberschrift.}
  \label{tab:a1-ergebnisse}
  \begin{tabular}{lrr}
    \toprule
    Variante & Wert A & Wert B \\
    \midrule
    Eins     & 1,2    & 3,4    \\
    Zwei     & 5,6    & 7,8    \\
    \bottomrule
  \end{tabular}
\end{table}

\section{Kennzeichnung von KI-Nutzung}

% ---------------------------------------------------------------------
\begin{beitraege}
  Recherche und Literaturauswahl                 & Max Mustermann \\
\end{beitraege}
